% ============================================================================
% comandos_eddie.tex — paquetes y comandos que usan los capítulos de la tesis.
% En tu proyecto de Overleaf: agrega % ============================================================================
% comandos_eddie.tex — paquetes y comandos que usan los capítulos de la tesis.
% En tu proyecto de Overleaf: agrega % ============================================================================
% comandos_eddie.tex — paquetes y comandos que usan los capítulos de la tesis.
% En tu proyecto de Overleaf: agrega % ============================================================================
% comandos_eddie.tex — paquetes y comandos que usan los capítulos de la tesis.
% En tu proyecto de Overleaf: agrega \input{comandos_eddie} al preámbulo de tu
% plantilla (después de sus propios paquetes). Si la plantilla ya carga alguno
% de estos paquetes no pasa nada: LaTeX no los carga dos veces.
% ============================================================================
\usepackage{graphicx}
\usepackage{booktabs}
\usepackage{tabularx}
\usepackage{longtable}
\usepackage{array}
\usepackage{multirow}
\usepackage{amsmath}
\usepackage{amssymb}
\usepackage[table]{xcolor}
\usepackage{float}
\usepackage{tikz}
\usetikzlibrary{arrows.meta,positioning,fit,backgrounds,shapes.geometric}
\usepackage[ruled,vlined,linesnumbered,spanish,onelanguage]{algorithm2e}
\usepackage{tcolorbox}
\usepackage{xurl}  % \nolinkurl corta nombres de archivo largos en cualquier punto

\definecolor{eddietinta}{HTML}{1B2440}
\definecolor{eddieacento}{HTML}{0B7A87}
\definecolor{eddiepend}{HTML}{B26A00}
\definecolor{eddiefondo}{HTML}{EEF1F6}

% Pie de tablas y figuras: «Fuente: …» (igual que en la plantilla del borrador)
\providecommand{\fuente}[1]{\par\vspace{2pt}{\footnotesize\centering\textit{Fuente:} #1\par}}

% Resultado que todavía no existe (se mide en el laboratorio). Se ve en el PDF para no olvidarlo.
\newcommand{\pendiente}[1]{{\color{eddiepend}\footnotesize\textsf{[Pendiente:} #1\textsf{]}}}

% Etiqueta del tipo de medición de cada prototipo
\newcommand{\tipomed}[1]{\textsf{\footnotesize[#1]}}

% Ficha de validación de un prototipo (misma estructura del borrador)
% #1 título, #2 etiqueta, #3 objetivo, #4 librerías/hardware, #5 contrato, #6 requisitos, #7 resultados
\newcommand{\ficha}[7]{%
\begin{table}[H]
\centering
\caption{Ficha de validación del prototipo: #1.}
\label{#2}
\small
\begin{tabularx}{\textwidth}{>{\bfseries}p{3.2cm} X}
\toprule
Objetivo & #3 \\ \midrule
Librerías / Hardware & #4 \\ \midrule
Contrato destino & #5 \\ \midrule
Requisitos asociados & #6 \\ \midrule
Resultados y hallazgos & #7 \\ \bottomrule
\end{tabularx}
\fuente{Elaboración propia, con \texttt{mediciones/medir\_prototipos.py} (P-00).}
\end{table}}

% Resultados numéricos generados a partir de las evidencias (ver datos/actualizar_datos.py)
\input{datos/prototipos_datos}
\input{datos/avances_datos}

% Figura que está en tu Overleaf pero no en este paquete: si el archivo existe se
% incluye; si no, se dibuja un recuadro con el nombre que debe tener.
\newcommand{\figuraopcional}[2]{%
  \IfFileExists{#1}{\includegraphics[width=#2]{#1}}{%
    \fbox{\parbox[c][4cm][c]{#2}{\centering\footnotesize Subir a Overleaf:\\ \texttt{\detokenize{#1}}}}}}

% Nombres largos de código (con guion bajo) pueden cortarse al final de línea.
\let\eddieunderscore\_
\renewcommand{\_}{\eddieunderscore\allowbreak}
\setlength{\emergencystretch}{3em}
 al preámbulo de tu
% plantilla (después de sus propios paquetes). Si la plantilla ya carga alguno
% de estos paquetes no pasa nada: LaTeX no los carga dos veces.
% ============================================================================
\usepackage{graphicx}
\usepackage{booktabs}
\usepackage{tabularx}
\usepackage{longtable}
\usepackage{array}
\usepackage{multirow}
\usepackage{amsmath}
\usepackage{amssymb}
\usepackage[table]{xcolor}
\usepackage{float}
\usepackage{tikz}
\usetikzlibrary{arrows.meta,positioning,fit,backgrounds,shapes.geometric}
\usepackage[ruled,vlined,linesnumbered,spanish,onelanguage]{algorithm2e}
\usepackage{tcolorbox}
\usepackage{xurl}  % \nolinkurl corta nombres de archivo largos en cualquier punto

\definecolor{eddietinta}{HTML}{1B2440}
\definecolor{eddieacento}{HTML}{0B7A87}
\definecolor{eddiepend}{HTML}{B26A00}
\definecolor{eddiefondo}{HTML}{EEF1F6}

% Pie de tablas y figuras: «Fuente: …» (igual que en la plantilla del borrador)
\providecommand{\fuente}[1]{\par\vspace{2pt}{\footnotesize\centering\textit{Fuente:} #1\par}}

% Resultado que todavía no existe (se mide en el laboratorio). Se ve en el PDF para no olvidarlo.
\newcommand{\pendiente}[1]{{\color{eddiepend}\footnotesize\textsf{[Pendiente:} #1\textsf{]}}}

% Etiqueta del tipo de medición de cada prototipo
\newcommand{\tipomed}[1]{\textsf{\footnotesize[#1]}}

% Ficha de validación de un prototipo (misma estructura del borrador)
% #1 título, #2 etiqueta, #3 objetivo, #4 librerías/hardware, #5 contrato, #6 requisitos, #7 resultados
\newcommand{\ficha}[7]{%
\begin{table}[H]
\centering
\caption{Ficha de validación del prototipo: #1.}
\label{#2}
\small
\begin{tabularx}{\textwidth}{>{\bfseries}p{3.2cm} X}
\toprule
Objetivo & #3 \\ \midrule
Librerías / Hardware & #4 \\ \midrule
Contrato destino & #5 \\ \midrule
Requisitos asociados & #6 \\ \midrule
Resultados y hallazgos & #7 \\ \bottomrule
\end{tabularx}
\fuente{Elaboración propia, con \texttt{mediciones/medir\_prototipos.py} (P-00).}
\end{table}}

% Resultados numéricos generados a partir de las evidencias (ver datos/actualizar_datos.py)
% Generado por mediciones/medir_prototipos.py — no editar a mano; volver a ejecutar el script.
\newcommand{\MedicionFecha}{2026-11-08}
\newcommand{\MedicionSistema}{Linux 6.18.44-fc-v310}
\newcommand{\MedicionPDF}{\texttt{borrador\_tesis.pdf}}
\newcommand{\FUnoTipo}{laboratorio}
\newcommand{\FDosPaginas}{30}
\newcommand{\FDosEficaciaPct}{25{,}8}
\newcommand{\FDosEficaciaSd}{2{,}7}
\newcommand{\FDosFTintaPct}{32{,}2}
\newcommand{\FDosEficaciaAdaptativoPct}{96{,}8}
\newcommand{\FDosTiempoMs}{0{,}67}
\newcommand{\FDosTiempoSdMs}{0{,}07}
\newcommand{\FDosTipo}{sintético}
\newcommand{\FTresBloques}{60}
\newcommand{\FTresCerLimpioPct}{1{,}46}
\newcommand{\FTresWerLimpioPct}{6{,}86}
\newcommand{\FTresCerDegradadoPct}{2{,}53}
\newcommand{\FTresWerDegradadoPct}{9{,}27}
\newcommand{\FTresCerSinPreprocesoPct}{0{,}42}
\newcommand{\FTresWerSinPreprocesoPct}{0{,}71}
\newcommand{\FTresCaracteres}{16353}
\newcommand{\FTresPalabras}{2405}
\newcommand{\FTresTiempoMs}{381{,}1}
\newcommand{\FTresTiempoSdMs}{48{,}2}
\newcommand{\FTresTipo}{pdf}
\newcommand{\FCuatroEscenas}{80}
\newcommand{\FCuatroExitoPct}{97{,}5}
\newcommand{\FCuatroRmsePx}{2{,}26}
\newcommand{\FCuatroRmseSdPx}{0{,}14}
\newcommand{\FCuatroTipo}{sintético}
\newcommand{\FCincoCuadros}{600}
\newcommand{\FCincoDeteccionPct}{100}
\newcommand{\FCincoFalsasPor}{0}
\newcommand{\FCincoDeteccionRojoPct}{0}
\newcommand{\FCincoFalsasRojoPor}{99{,}7}
\newcommand{\FCincoErrorPx}{0{,}73}
\newcommand{\FCincoErrorSdPx}{0{,}3}
\newcommand{\FCincoTiempoMs}{1{,}93}
\newcommand{\FCincoTipo}{sintético}
\newcommand{\FSeisCuadros}{200}
\newcommand{\FSeisIouPct}{96{,}5}
\newcommand{\FSeisPuntaOkPct}{96{,}5}
\newcommand{\FSeisErrorPuntaPx}{1{,}2}
\newcommand{\FSeisTiempoMs}{2{,}17}
\newcommand{\FSeisTipo}{sintético}
\newcommand{\FSieteEnsayos}{300}
\newcommand{\FSieteErrorPx}{1{,}53}
\newcommand{\FSieteErrorSdPx}{0{,}51}
\newcommand{\FSietePPx}{2{,}46}
\newcommand{\FSieteTipo}{sintético}
\newcommand{\FNueveTipo}{laboratorio}
\newcommand{\FDiezMuestras}{3000}
\newcommand{\FDiezRecibidas}{3000}
\newcommand{\FDiezPerdidaPct}{0}
\newcommand{\FDiezFrecuenciaHz}{60}
\newcommand{\FDiezJitterMs}{0{,}45}
\newcommand{\FDiezTipo}{sintético}
\newcommand{\FOnceModulos}{8}
\newcommand{\FOnceCargas}{240}
\newcommand{\FOnceFallas}{0}
\newcommand{\FOnceTiempoMs}{0{,}36}
\newcommand{\FOnceTiempoSdMs}{0{,}79}
\newcommand{\FOnceTipo}{código}
\newcommand{\FDoceUmbralS}{0{,}6}
\newcommand{\FDoceRadioPx}{45}
\newcommand{\FDoceVoluntariasPct}{100}
\newcommand{\FDoceAccidentalesPorMin}{3{,}5}
\newcommand{\FDoceMiradasVoluntarias}{30}
\newcommand{\FDoceVoluntariasRuidoBajoPct}{100}
\newcommand{\FDoceAccidentalesRuidoBajo}{7{,}9}
\newcommand{\FDoceVoluntariasRuidoAltoPct}{50}
\newcommand{\FDoceAccidentalesRuidoAlto}{0{,}3}
\newcommand{\FDoceBarrido}{0{,}4 & 25 & 33{,}3 & 1 \\
0{,}4 & 35 & 76{,}7 & 10{,}5 \\
0{,}4 & 45 & 100 & 29{,}3 \\
0{,}5 & 25 & 26{,}7 & 0 \\
0{,}5 & 35 & 76{,}7 & 1{,}4 \\
0{,}5 & 45 & 100 & 9{,}6 \\
0{,}6 & 25 & 13{,}3 & 0 \\
0{,}6 & 35 & 66{,}7 & 0{,}4 \\
0{,}6 & 45 & 100 & 3{,}5 \\
0{,}8 & 25 & 3{,}3 & 0 \\
0{,}8 & 35 & 50 & 0 \\
0{,}8 & 45 & 93{,}3 & 0{,}3}
\newcommand{\FDoceTipo}{sintético}
\newcommand{\FTreceTiempoMs}{0{,}82}
\newcommand{\FTreceTiempoSdMs}{0{,}06}
\newcommand{\FTreceFpsMax}{1224{,}6}
\newcommand{\FTreceTipo}{sintético}
\newcommand{\FCatorcePaginas}{64}
\newcommand{\FCatorceConsultas}{1920}
\newcommand{\FCatorceCoincidencias}{1836}
\newcommand{\FCatorceCoherenciaPct}{82}
\newcommand{\FCatorceTiempoMs}{1{,}03}
\newcommand{\FCatorceTiempoSdMs}{0{,}62}
\newcommand{\FCatorceTipo}{pdf}
\newcommand{\FQuinceAnotaciones}{48}
\newcommand{\FQuincePersistidas}{48}
\newcommand{\FQuinceTiempoMs}{13{,}8}
\newcommand{\FQuinceTiempoSdMs}{1{,}7}
\newcommand{\FQuinceTipo}{pdf}
\newcommand{\FDieciseisNotas}{40}
\newcommand{\FDieciseisDetectadasPct}{100}
\newcommand{\FDieciseisErrorPt}{0{,}68}
\newcommand{\FDieciseisErrorSdPt}{0}
\newcommand{\FDieciseisTipo}{sintético}
\newcommand{\FDiecisieteFiguras}{60}
\newcommand{\FDiecisieteExactitudPct}{80}
\newcommand{\FDiecisieteIouPct}{90{,}1}
\newcommand{\FDiecisieteIouSd}{2{,}8}
\newcommand{\FDiecisieteTipo}{sintético}
\newcommand{\FVeinticuatroTipo}{laboratorio}
\newcommand{\FVeinticincoTipo}{laboratorio}
\newcommand{\FVeintisieteAlfa}{0{,}3}
\newcommand{\FVeintisieteTiempoMs}{0{,}58}
\newcommand{\FVeintisieteTiempoSdMs}{0{,}07}
\newcommand{\FVeintisieteTipo}{sintético}

% Generado por datos/actualizar_datos.py a partir de las evidencias de los avances. No editar.
\newcommand{\PUnoPruebas}{34}
\newcommand{\PUnoXfail}{0}
\newcommand{\PUnoCobertura}{81}
\newcommand{\PUnoCicloMedia}{0{,}75}
\newcommand{\PUnoCicloPnoventaycinco}{0{,}98}
\newcommand{\PUnoCicloMax}{2{,}5}
\newcommand{\PUnoFecha}{2026-11-08}
\newcommand{\PDosPruebas}{41}
\newcommand{\PDosXfail}{1}
\newcommand{\PDosCobertura}{73}
\newcommand{\PDosCicloMedia}{0{,}15}
\newcommand{\PDosCicloPnoventaycinco}{0{,}2}
\newcommand{\PDosCicloMax}{2{,}49}
\newcommand{\PDosFecha}{2026-11-08}
\newcommand{\PTresPruebas}{65}
\newcommand{\PTresXfail}{0}
\newcommand{\PTresCobertura}{67}
\newcommand{\PTresCicloMedia}{0{,}15}
\newcommand{\PTresCicloPnoventaycinco}{0{,}2}
\newcommand{\PTresCicloMax}{0{,}58}
\newcommand{\PTresFecha}{2026-11-08}
\newcommand{\PCuatroPruebas}{79}
\newcommand{\PCuatroXfail}{1}
\newcommand{\PCuatroCobertura}{72}
\newcommand{\PCuatroCicloMedia}{0{,}6}
\newcommand{\PCuatroCicloPnoventaycinco}{1{,}5}
\newcommand{\PCuatroCicloMax}{1{,}76}
\newcommand{\PCuatroFecha}{2026-11-08}
\newcommand{\PCincoPruebas}{94}
\newcommand{\PCincoXfail}{1}
\newcommand{\PCincoCobertura}{74}
\newcommand{\PCincoCicloMedia}{0{,}58}
\newcommand{\PCincoCicloPnoventaycinco}{1{,}41}
\newcommand{\PCincoCicloMax}{3{,}53}
\newcommand{\PCincoFecha}{2026-11-08}
\newcommand{\PSeisPruebas}{100}
\newcommand{\PSeisXfail}{2}
\newcommand{\PSeisCobertura}{81}
\newcommand{\PSeisCicloMedia}{0{,}66}
\newcommand{\PSeisCicloPnoventaycinco}{1{,}49}
\newcommand{\PSeisCicloMax}{10{,}28}
\newcommand{\PSeisFecha}{2026-11-08}
\newcommand{\PSietePruebas}{107}
\newcommand{\PSieteXfail}{0}
\newcommand{\PSieteCobertura}{81}
\newcommand{\PSieteCicloMedia}{0{,}59}
\newcommand{\PSieteCicloPnoventaycinco}{1{,}48}
\newcommand{\PSieteCicloMax}{2{,}31}
\newcommand{\PSieteFecha}{2026-11-08}
\newcommand{\PSieteCiclos}{500}
\newcommand{\PSieteTablaLatencias}{TextCamera & 0{,}117 & 0{,}182 & 0{,}76 \\
GestureCamera & 0{,}264 & 0{,}369 & 2{,}12 \\
ImageProcessor & 0{,}006 & 0{,}008 & 0{,}057 \\
ColorPenRecognition & 0{,}122 & 0{,}932 & 1{,}189 \\
HandSkinRecognition & 0{,}003 & 0{,}005 & 0{,}046 \\
MouseRecognition & 0{,}002 & 0{,}003 & 0{,}013 \\
EyeTrackerMock & 0{,}01 & 0{,}015 & 0{,}16 \\
SpeechProvider & 0{,}003 & 0{,}004 & 0{,}079 \\
GazeSession & 0{,}005 & 0{,}006 & 0{,}13 \\
ConsistencyProvider & 0{,}009 & 0{,}011 & 0{,}074}
\newcommand{\PSieteEntorno}{Linux 6.18.44-fc-v310, Python 3.13.18}


% Figura que está en tu Overleaf pero no en este paquete: si el archivo existe se
% incluye; si no, se dibuja un recuadro con el nombre que debe tener.
\newcommand{\figuraopcional}[2]{%
  \IfFileExists{#1}{\includegraphics[width=#2]{#1}}{%
    \fbox{\parbox[c][4cm][c]{#2}{\centering\footnotesize Subir a Overleaf:\\ \texttt{\detokenize{#1}}}}}}

% Nombres largos de código (con guion bajo) pueden cortarse al final de línea.
\let\eddieunderscore\_
\renewcommand{\_}{\eddieunderscore\allowbreak}
\setlength{\emergencystretch}{3em}
 al preámbulo de tu
% plantilla (después de sus propios paquetes). Si la plantilla ya carga alguno
% de estos paquetes no pasa nada: LaTeX no los carga dos veces.
% ============================================================================
\usepackage{graphicx}
\usepackage{booktabs}
\usepackage{tabularx}
\usepackage{longtable}
\usepackage{array}
\usepackage{multirow}
\usepackage{amsmath}
\usepackage{amssymb}
\usepackage[table]{xcolor}
\usepackage{float}
\usepackage{tikz}
\usetikzlibrary{arrows.meta,positioning,fit,backgrounds,shapes.geometric}
\usepackage[ruled,vlined,linesnumbered,spanish,onelanguage]{algorithm2e}
\usepackage{tcolorbox}
\usepackage{xurl}  % \nolinkurl corta nombres de archivo largos en cualquier punto

\definecolor{eddietinta}{HTML}{1B2440}
\definecolor{eddieacento}{HTML}{0B7A87}
\definecolor{eddiepend}{HTML}{B26A00}
\definecolor{eddiefondo}{HTML}{EEF1F6}

% Pie de tablas y figuras: «Fuente: …» (igual que en la plantilla del borrador)
\providecommand{\fuente}[1]{\par\vspace{2pt}{\footnotesize\centering\textit{Fuente:} #1\par}}

% Resultado que todavía no existe (se mide en el laboratorio). Se ve en el PDF para no olvidarlo.
\newcommand{\pendiente}[1]{{\color{eddiepend}\footnotesize\textsf{[Pendiente:} #1\textsf{]}}}

% Etiqueta del tipo de medición de cada prototipo
\newcommand{\tipomed}[1]{\textsf{\footnotesize[#1]}}

% Ficha de validación de un prototipo (misma estructura del borrador)
% #1 título, #2 etiqueta, #3 objetivo, #4 librerías/hardware, #5 contrato, #6 requisitos, #7 resultados
\newcommand{\ficha}[7]{%
\begin{table}[H]
\centering
\caption{Ficha de validación del prototipo: #1.}
\label{#2}
\small
\begin{tabularx}{\textwidth}{>{\bfseries}p{3.2cm} X}
\toprule
Objetivo & #3 \\ \midrule
Librerías / Hardware & #4 \\ \midrule
Contrato destino & #5 \\ \midrule
Requisitos asociados & #6 \\ \midrule
Resultados y hallazgos & #7 \\ \bottomrule
\end{tabularx}
\fuente{Elaboración propia, con \texttt{mediciones/medir\_prototipos.py} (P-00).}
\end{table}}

% Resultados numéricos generados a partir de las evidencias (ver datos/actualizar_datos.py)
% Generado por mediciones/medir_prototipos.py — no editar a mano; volver a ejecutar el script.
\newcommand{\MedicionFecha}{2026-11-08}
\newcommand{\MedicionSistema}{Linux 6.18.44-fc-v310}
\newcommand{\MedicionPDF}{\texttt{borrador\_tesis.pdf}}
\newcommand{\FUnoTipo}{laboratorio}
\newcommand{\FDosPaginas}{30}
\newcommand{\FDosEficaciaPct}{25{,}8}
\newcommand{\FDosEficaciaSd}{2{,}7}
\newcommand{\FDosFTintaPct}{32{,}2}
\newcommand{\FDosEficaciaAdaptativoPct}{96{,}8}
\newcommand{\FDosTiempoMs}{0{,}67}
\newcommand{\FDosTiempoSdMs}{0{,}07}
\newcommand{\FDosTipo}{sintético}
\newcommand{\FTresBloques}{60}
\newcommand{\FTresCerLimpioPct}{1{,}46}
\newcommand{\FTresWerLimpioPct}{6{,}86}
\newcommand{\FTresCerDegradadoPct}{2{,}53}
\newcommand{\FTresWerDegradadoPct}{9{,}27}
\newcommand{\FTresCerSinPreprocesoPct}{0{,}42}
\newcommand{\FTresWerSinPreprocesoPct}{0{,}71}
\newcommand{\FTresCaracteres}{16353}
\newcommand{\FTresPalabras}{2405}
\newcommand{\FTresTiempoMs}{381{,}1}
\newcommand{\FTresTiempoSdMs}{48{,}2}
\newcommand{\FTresTipo}{pdf}
\newcommand{\FCuatroEscenas}{80}
\newcommand{\FCuatroExitoPct}{97{,}5}
\newcommand{\FCuatroRmsePx}{2{,}26}
\newcommand{\FCuatroRmseSdPx}{0{,}14}
\newcommand{\FCuatroTipo}{sintético}
\newcommand{\FCincoCuadros}{600}
\newcommand{\FCincoDeteccionPct}{100}
\newcommand{\FCincoFalsasPor}{0}
\newcommand{\FCincoDeteccionRojoPct}{0}
\newcommand{\FCincoFalsasRojoPor}{99{,}7}
\newcommand{\FCincoErrorPx}{0{,}73}
\newcommand{\FCincoErrorSdPx}{0{,}3}
\newcommand{\FCincoTiempoMs}{1{,}93}
\newcommand{\FCincoTipo}{sintético}
\newcommand{\FSeisCuadros}{200}
\newcommand{\FSeisIouPct}{96{,}5}
\newcommand{\FSeisPuntaOkPct}{96{,}5}
\newcommand{\FSeisErrorPuntaPx}{1{,}2}
\newcommand{\FSeisTiempoMs}{2{,}17}
\newcommand{\FSeisTipo}{sintético}
\newcommand{\FSieteEnsayos}{300}
\newcommand{\FSieteErrorPx}{1{,}53}
\newcommand{\FSieteErrorSdPx}{0{,}51}
\newcommand{\FSietePPx}{2{,}46}
\newcommand{\FSieteTipo}{sintético}
\newcommand{\FNueveTipo}{laboratorio}
\newcommand{\FDiezMuestras}{3000}
\newcommand{\FDiezRecibidas}{3000}
\newcommand{\FDiezPerdidaPct}{0}
\newcommand{\FDiezFrecuenciaHz}{60}
\newcommand{\FDiezJitterMs}{0{,}45}
\newcommand{\FDiezTipo}{sintético}
\newcommand{\FOnceModulos}{8}
\newcommand{\FOnceCargas}{240}
\newcommand{\FOnceFallas}{0}
\newcommand{\FOnceTiempoMs}{0{,}36}
\newcommand{\FOnceTiempoSdMs}{0{,}79}
\newcommand{\FOnceTipo}{código}
\newcommand{\FDoceUmbralS}{0{,}6}
\newcommand{\FDoceRadioPx}{45}
\newcommand{\FDoceVoluntariasPct}{100}
\newcommand{\FDoceAccidentalesPorMin}{3{,}5}
\newcommand{\FDoceMiradasVoluntarias}{30}
\newcommand{\FDoceVoluntariasRuidoBajoPct}{100}
\newcommand{\FDoceAccidentalesRuidoBajo}{7{,}9}
\newcommand{\FDoceVoluntariasRuidoAltoPct}{50}
\newcommand{\FDoceAccidentalesRuidoAlto}{0{,}3}
\newcommand{\FDoceBarrido}{0{,}4 & 25 & 33{,}3 & 1 \\
0{,}4 & 35 & 76{,}7 & 10{,}5 \\
0{,}4 & 45 & 100 & 29{,}3 \\
0{,}5 & 25 & 26{,}7 & 0 \\
0{,}5 & 35 & 76{,}7 & 1{,}4 \\
0{,}5 & 45 & 100 & 9{,}6 \\
0{,}6 & 25 & 13{,}3 & 0 \\
0{,}6 & 35 & 66{,}7 & 0{,}4 \\
0{,}6 & 45 & 100 & 3{,}5 \\
0{,}8 & 25 & 3{,}3 & 0 \\
0{,}8 & 35 & 50 & 0 \\
0{,}8 & 45 & 93{,}3 & 0{,}3}
\newcommand{\FDoceTipo}{sintético}
\newcommand{\FTreceTiempoMs}{0{,}82}
\newcommand{\FTreceTiempoSdMs}{0{,}06}
\newcommand{\FTreceFpsMax}{1224{,}6}
\newcommand{\FTreceTipo}{sintético}
\newcommand{\FCatorcePaginas}{64}
\newcommand{\FCatorceConsultas}{1920}
\newcommand{\FCatorceCoincidencias}{1836}
\newcommand{\FCatorceCoherenciaPct}{82}
\newcommand{\FCatorceTiempoMs}{1{,}03}
\newcommand{\FCatorceTiempoSdMs}{0{,}62}
\newcommand{\FCatorceTipo}{pdf}
\newcommand{\FQuinceAnotaciones}{48}
\newcommand{\FQuincePersistidas}{48}
\newcommand{\FQuinceTiempoMs}{13{,}8}
\newcommand{\FQuinceTiempoSdMs}{1{,}7}
\newcommand{\FQuinceTipo}{pdf}
\newcommand{\FDieciseisNotas}{40}
\newcommand{\FDieciseisDetectadasPct}{100}
\newcommand{\FDieciseisErrorPt}{0{,}68}
\newcommand{\FDieciseisErrorSdPt}{0}
\newcommand{\FDieciseisTipo}{sintético}
\newcommand{\FDiecisieteFiguras}{60}
\newcommand{\FDiecisieteExactitudPct}{80}
\newcommand{\FDiecisieteIouPct}{90{,}1}
\newcommand{\FDiecisieteIouSd}{2{,}8}
\newcommand{\FDiecisieteTipo}{sintético}
\newcommand{\FVeinticuatroTipo}{laboratorio}
\newcommand{\FVeinticincoTipo}{laboratorio}
\newcommand{\FVeintisieteAlfa}{0{,}3}
\newcommand{\FVeintisieteTiempoMs}{0{,}58}
\newcommand{\FVeintisieteTiempoSdMs}{0{,}07}
\newcommand{\FVeintisieteTipo}{sintético}

% Generado por datos/actualizar_datos.py a partir de las evidencias de los avances. No editar.
\newcommand{\PUnoPruebas}{34}
\newcommand{\PUnoXfail}{0}
\newcommand{\PUnoCobertura}{81}
\newcommand{\PUnoCicloMedia}{0{,}75}
\newcommand{\PUnoCicloPnoventaycinco}{0{,}98}
\newcommand{\PUnoCicloMax}{2{,}5}
\newcommand{\PUnoFecha}{2026-11-08}
\newcommand{\PDosPruebas}{41}
\newcommand{\PDosXfail}{1}
\newcommand{\PDosCobertura}{73}
\newcommand{\PDosCicloMedia}{0{,}15}
\newcommand{\PDosCicloPnoventaycinco}{0{,}2}
\newcommand{\PDosCicloMax}{2{,}49}
\newcommand{\PDosFecha}{2026-11-08}
\newcommand{\PTresPruebas}{65}
\newcommand{\PTresXfail}{0}
\newcommand{\PTresCobertura}{67}
\newcommand{\PTresCicloMedia}{0{,}15}
\newcommand{\PTresCicloPnoventaycinco}{0{,}2}
\newcommand{\PTresCicloMax}{0{,}58}
\newcommand{\PTresFecha}{2026-11-08}
\newcommand{\PCuatroPruebas}{79}
\newcommand{\PCuatroXfail}{1}
\newcommand{\PCuatroCobertura}{72}
\newcommand{\PCuatroCicloMedia}{0{,}6}
\newcommand{\PCuatroCicloPnoventaycinco}{1{,}5}
\newcommand{\PCuatroCicloMax}{1{,}76}
\newcommand{\PCuatroFecha}{2026-11-08}
\newcommand{\PCincoPruebas}{94}
\newcommand{\PCincoXfail}{1}
\newcommand{\PCincoCobertura}{74}
\newcommand{\PCincoCicloMedia}{0{,}58}
\newcommand{\PCincoCicloPnoventaycinco}{1{,}41}
\newcommand{\PCincoCicloMax}{3{,}53}
\newcommand{\PCincoFecha}{2026-11-08}
\newcommand{\PSeisPruebas}{100}
\newcommand{\PSeisXfail}{2}
\newcommand{\PSeisCobertura}{81}
\newcommand{\PSeisCicloMedia}{0{,}66}
\newcommand{\PSeisCicloPnoventaycinco}{1{,}49}
\newcommand{\PSeisCicloMax}{10{,}28}
\newcommand{\PSeisFecha}{2026-11-08}
\newcommand{\PSietePruebas}{107}
\newcommand{\PSieteXfail}{0}
\newcommand{\PSieteCobertura}{81}
\newcommand{\PSieteCicloMedia}{0{,}59}
\newcommand{\PSieteCicloPnoventaycinco}{1{,}48}
\newcommand{\PSieteCicloMax}{2{,}31}
\newcommand{\PSieteFecha}{2026-11-08}
\newcommand{\PSieteCiclos}{500}
\newcommand{\PSieteTablaLatencias}{TextCamera & 0{,}117 & 0{,}182 & 0{,}76 \\
GestureCamera & 0{,}264 & 0{,}369 & 2{,}12 \\
ImageProcessor & 0{,}006 & 0{,}008 & 0{,}057 \\
ColorPenRecognition & 0{,}122 & 0{,}932 & 1{,}189 \\
HandSkinRecognition & 0{,}003 & 0{,}005 & 0{,}046 \\
MouseRecognition & 0{,}002 & 0{,}003 & 0{,}013 \\
EyeTrackerMock & 0{,}01 & 0{,}015 & 0{,}16 \\
SpeechProvider & 0{,}003 & 0{,}004 & 0{,}079 \\
GazeSession & 0{,}005 & 0{,}006 & 0{,}13 \\
ConsistencyProvider & 0{,}009 & 0{,}011 & 0{,}074}
\newcommand{\PSieteEntorno}{Linux 6.18.44-fc-v310, Python 3.13.18}


% Figura que está en tu Overleaf pero no en este paquete: si el archivo existe se
% incluye; si no, se dibuja un recuadro con el nombre que debe tener.
\newcommand{\figuraopcional}[2]{%
  \IfFileExists{#1}{\includegraphics[width=#2]{#1}}{%
    \fbox{\parbox[c][4cm][c]{#2}{\centering\footnotesize Subir a Overleaf:\\ \texttt{\detokenize{#1}}}}}}

% Nombres largos de código (con guion bajo) pueden cortarse al final de línea.
\let\eddieunderscore\_
\renewcommand{\_}{\eddieunderscore\allowbreak}
\setlength{\emergencystretch}{3em}
 al preámbulo de tu
% plantilla (después de sus propios paquetes). Si la plantilla ya carga alguno
% de estos paquetes no pasa nada: LaTeX no los carga dos veces.
% ============================================================================
\usepackage{graphicx}
\usepackage{booktabs}
\usepackage{tabularx}
\usepackage{longtable}
\usepackage{array}
\usepackage{multirow}
\usepackage{amsmath}
\usepackage{amssymb}
\usepackage[table]{xcolor}
\usepackage{float}
\usepackage{tikz}
\usetikzlibrary{arrows.meta,positioning,fit,backgrounds,shapes.geometric}
\usepackage[ruled,vlined,linesnumbered,spanish,onelanguage]{algorithm2e}
\usepackage{tcolorbox}
\usepackage{xurl}  % \nolinkurl corta nombres de archivo largos en cualquier punto

\definecolor{eddietinta}{HTML}{1B2440}
\definecolor{eddieacento}{HTML}{0B7A87}
\definecolor{eddiepend}{HTML}{B26A00}
\definecolor{eddiefondo}{HTML}{EEF1F6}

% Pie de tablas y figuras: «Fuente: …» (igual que en la plantilla del borrador)
\providecommand{\fuente}[1]{\par\vspace{2pt}{\footnotesize\centering\textit{Fuente:} #1\par}}

% Resultado que todavía no existe (se mide en el laboratorio). Se ve en el PDF para no olvidarlo.
\newcommand{\pendiente}[1]{{\color{eddiepend}\footnotesize\textsf{[Pendiente:} #1\textsf{]}}}

% Etiqueta del tipo de medición de cada prototipo
\newcommand{\tipomed}[1]{\textsf{\footnotesize[#1]}}

% Ficha de validación de un prototipo (misma estructura del borrador)
% #1 título, #2 etiqueta, #3 objetivo, #4 librerías/hardware, #5 contrato, #6 requisitos, #7 resultados
\newcommand{\ficha}[7]{%
\begin{table}[H]
\centering
\caption{Ficha de validación del prototipo: #1.}
\label{#2}
\small
\begin{tabularx}{\textwidth}{>{\bfseries}p{3.2cm} X}
\toprule
Objetivo & #3 \\ \midrule
Librerías / Hardware & #4 \\ \midrule
Contrato destino & #5 \\ \midrule
Requisitos asociados & #6 \\ \midrule
Resultados y hallazgos & #7 \\ \bottomrule
\end{tabularx}
\fuente{Elaboración propia, con \texttt{mediciones/medir\_prototipos.py} (P-00).}
\end{table}}

% Resultados numéricos generados a partir de las evidencias (ver datos/actualizar_datos.py)
% Generado por mediciones/medir_prototipos.py — no editar a mano; volver a ejecutar el script.
\newcommand{\MedicionFecha}{2026-11-08}
\newcommand{\MedicionSistema}{Linux 6.18.44-fc-v310}
\newcommand{\MedicionPDF}{\texttt{borrador\_tesis.pdf}}
\newcommand{\FUnoTipo}{laboratorio}
\newcommand{\FDosPaginas}{30}
\newcommand{\FDosEficaciaPct}{25{,}8}
\newcommand{\FDosEficaciaSd}{2{,}7}
\newcommand{\FDosFTintaPct}{32{,}2}
\newcommand{\FDosEficaciaAdaptativoPct}{96{,}8}
\newcommand{\FDosTiempoMs}{0{,}67}
\newcommand{\FDosTiempoSdMs}{0{,}07}
\newcommand{\FDosTipo}{sintético}
\newcommand{\FTresBloques}{60}
\newcommand{\FTresCerLimpioPct}{1{,}46}
\newcommand{\FTresWerLimpioPct}{6{,}86}
\newcommand{\FTresCerDegradadoPct}{2{,}53}
\newcommand{\FTresWerDegradadoPct}{9{,}27}
\newcommand{\FTresCerSinPreprocesoPct}{0{,}42}
\newcommand{\FTresWerSinPreprocesoPct}{0{,}71}
\newcommand{\FTresCaracteres}{16353}
\newcommand{\FTresPalabras}{2405}
\newcommand{\FTresTiempoMs}{381{,}1}
\newcommand{\FTresTiempoSdMs}{48{,}2}
\newcommand{\FTresTipo}{pdf}
\newcommand{\FCuatroEscenas}{80}
\newcommand{\FCuatroExitoPct}{97{,}5}
\newcommand{\FCuatroRmsePx}{2{,}26}
\newcommand{\FCuatroRmseSdPx}{0{,}14}
\newcommand{\FCuatroTipo}{sintético}
\newcommand{\FCincoCuadros}{600}
\newcommand{\FCincoDeteccionPct}{100}
\newcommand{\FCincoFalsasPor}{0}
\newcommand{\FCincoDeteccionRojoPct}{0}
\newcommand{\FCincoFalsasRojoPor}{99{,}7}
\newcommand{\FCincoErrorPx}{0{,}73}
\newcommand{\FCincoErrorSdPx}{0{,}3}
\newcommand{\FCincoTiempoMs}{1{,}93}
\newcommand{\FCincoTipo}{sintético}
\newcommand{\FSeisCuadros}{200}
\newcommand{\FSeisIouPct}{96{,}5}
\newcommand{\FSeisPuntaOkPct}{96{,}5}
\newcommand{\FSeisErrorPuntaPx}{1{,}2}
\newcommand{\FSeisTiempoMs}{2{,}17}
\newcommand{\FSeisTipo}{sintético}
\newcommand{\FSieteEnsayos}{300}
\newcommand{\FSieteErrorPx}{1{,}53}
\newcommand{\FSieteErrorSdPx}{0{,}51}
\newcommand{\FSietePPx}{2{,}46}
\newcommand{\FSieteTipo}{sintético}
\newcommand{\FNueveTipo}{laboratorio}
\newcommand{\FDiezMuestras}{3000}
\newcommand{\FDiezRecibidas}{3000}
\newcommand{\FDiezPerdidaPct}{0}
\newcommand{\FDiezFrecuenciaHz}{60}
\newcommand{\FDiezJitterMs}{0{,}45}
\newcommand{\FDiezTipo}{sintético}
\newcommand{\FOnceModulos}{8}
\newcommand{\FOnceCargas}{240}
\newcommand{\FOnceFallas}{0}
\newcommand{\FOnceTiempoMs}{0{,}36}
\newcommand{\FOnceTiempoSdMs}{0{,}79}
\newcommand{\FOnceTipo}{código}
\newcommand{\FDoceUmbralS}{0{,}6}
\newcommand{\FDoceRadioPx}{45}
\newcommand{\FDoceVoluntariasPct}{100}
\newcommand{\FDoceAccidentalesPorMin}{3{,}5}
\newcommand{\FDoceMiradasVoluntarias}{30}
\newcommand{\FDoceVoluntariasRuidoBajoPct}{100}
\newcommand{\FDoceAccidentalesRuidoBajo}{7{,}9}
\newcommand{\FDoceVoluntariasRuidoAltoPct}{50}
\newcommand{\FDoceAccidentalesRuidoAlto}{0{,}3}
\newcommand{\FDoceBarrido}{0{,}4 & 25 & 33{,}3 & 1 \\
0{,}4 & 35 & 76{,}7 & 10{,}5 \\
0{,}4 & 45 & 100 & 29{,}3 \\
0{,}5 & 25 & 26{,}7 & 0 \\
0{,}5 & 35 & 76{,}7 & 1{,}4 \\
0{,}5 & 45 & 100 & 9{,}6 \\
0{,}6 & 25 & 13{,}3 & 0 \\
0{,}6 & 35 & 66{,}7 & 0{,}4 \\
0{,}6 & 45 & 100 & 3{,}5 \\
0{,}8 & 25 & 3{,}3 & 0 \\
0{,}8 & 35 & 50 & 0 \\
0{,}8 & 45 & 93{,}3 & 0{,}3}
\newcommand{\FDoceTipo}{sintético}
\newcommand{\FTreceTiempoMs}{0{,}82}
\newcommand{\FTreceTiempoSdMs}{0{,}06}
\newcommand{\FTreceFpsMax}{1224{,}6}
\newcommand{\FTreceTipo}{sintético}
\newcommand{\FCatorcePaginas}{64}
\newcommand{\FCatorceConsultas}{1920}
\newcommand{\FCatorceCoincidencias}{1836}
\newcommand{\FCatorceCoherenciaPct}{82}
\newcommand{\FCatorceTiempoMs}{1{,}03}
\newcommand{\FCatorceTiempoSdMs}{0{,}62}
\newcommand{\FCatorceTipo}{pdf}
\newcommand{\FQuinceAnotaciones}{48}
\newcommand{\FQuincePersistidas}{48}
\newcommand{\FQuinceTiempoMs}{13{,}8}
\newcommand{\FQuinceTiempoSdMs}{1{,}7}
\newcommand{\FQuinceTipo}{pdf}
\newcommand{\FDieciseisNotas}{40}
\newcommand{\FDieciseisDetectadasPct}{100}
\newcommand{\FDieciseisErrorPt}{0{,}68}
\newcommand{\FDieciseisErrorSdPt}{0}
\newcommand{\FDieciseisTipo}{sintético}
\newcommand{\FDiecisieteFiguras}{60}
\newcommand{\FDiecisieteExactitudPct}{80}
\newcommand{\FDiecisieteIouPct}{90{,}1}
\newcommand{\FDiecisieteIouSd}{2{,}8}
\newcommand{\FDiecisieteTipo}{sintético}
\newcommand{\FVeinticuatroTipo}{laboratorio}
\newcommand{\FVeinticincoTipo}{laboratorio}
\newcommand{\FVeintisieteAlfa}{0{,}3}
\newcommand{\FVeintisieteTiempoMs}{0{,}58}
\newcommand{\FVeintisieteTiempoSdMs}{0{,}07}
\newcommand{\FVeintisieteTipo}{sintético}

% Generado por datos/actualizar_datos.py a partir de las evidencias de los avances. No editar.
\newcommand{\PUnoPruebas}{34}
\newcommand{\PUnoXfail}{0}
\newcommand{\PUnoCobertura}{81}
\newcommand{\PUnoCicloMedia}{0{,}75}
\newcommand{\PUnoCicloPnoventaycinco}{0{,}98}
\newcommand{\PUnoCicloMax}{2{,}5}
\newcommand{\PUnoFecha}{2026-11-08}
\newcommand{\PDosPruebas}{41}
\newcommand{\PDosXfail}{1}
\newcommand{\PDosCobertura}{73}
\newcommand{\PDosCicloMedia}{0{,}15}
\newcommand{\PDosCicloPnoventaycinco}{0{,}2}
\newcommand{\PDosCicloMax}{2{,}49}
\newcommand{\PDosFecha}{2026-11-08}
\newcommand{\PTresPruebas}{65}
\newcommand{\PTresXfail}{0}
\newcommand{\PTresCobertura}{67}
\newcommand{\PTresCicloMedia}{0{,}15}
\newcommand{\PTresCicloPnoventaycinco}{0{,}2}
\newcommand{\PTresCicloMax}{0{,}58}
\newcommand{\PTresFecha}{2026-11-08}
\newcommand{\PCuatroPruebas}{79}
\newcommand{\PCuatroXfail}{1}
\newcommand{\PCuatroCobertura}{72}
\newcommand{\PCuatroCicloMedia}{0{,}6}
\newcommand{\PCuatroCicloPnoventaycinco}{1{,}5}
\newcommand{\PCuatroCicloMax}{1{,}76}
\newcommand{\PCuatroFecha}{2026-11-08}
\newcommand{\PCincoPruebas}{94}
\newcommand{\PCincoXfail}{1}
\newcommand{\PCincoCobertura}{74}
\newcommand{\PCincoCicloMedia}{0{,}58}
\newcommand{\PCincoCicloPnoventaycinco}{1{,}41}
\newcommand{\PCincoCicloMax}{3{,}53}
\newcommand{\PCincoFecha}{2026-11-08}
\newcommand{\PSeisPruebas}{100}
\newcommand{\PSeisXfail}{2}
\newcommand{\PSeisCobertura}{81}
\newcommand{\PSeisCicloMedia}{0{,}66}
\newcommand{\PSeisCicloPnoventaycinco}{1{,}49}
\newcommand{\PSeisCicloMax}{10{,}28}
\newcommand{\PSeisFecha}{2026-11-08}
\newcommand{\PSietePruebas}{107}
\newcommand{\PSieteXfail}{0}
\newcommand{\PSieteCobertura}{81}
\newcommand{\PSieteCicloMedia}{0{,}59}
\newcommand{\PSieteCicloPnoventaycinco}{1{,}48}
\newcommand{\PSieteCicloMax}{2{,}31}
\newcommand{\PSieteFecha}{2026-11-08}
\newcommand{\PSieteCiclos}{500}
\newcommand{\PSieteTablaLatencias}{TextCamera & 0{,}117 & 0{,}182 & 0{,}76 \\
GestureCamera & 0{,}264 & 0{,}369 & 2{,}12 \\
ImageProcessor & 0{,}006 & 0{,}008 & 0{,}057 \\
ColorPenRecognition & 0{,}122 & 0{,}932 & 1{,}189 \\
HandSkinRecognition & 0{,}003 & 0{,}005 & 0{,}046 \\
MouseRecognition & 0{,}002 & 0{,}003 & 0{,}013 \\
EyeTrackerMock & 0{,}01 & 0{,}015 & 0{,}16 \\
SpeechProvider & 0{,}003 & 0{,}004 & 0{,}079 \\
GazeSession & 0{,}005 & 0{,}006 & 0{,}13 \\
ConsistencyProvider & 0{,}009 & 0{,}011 & 0{,}074}
\newcommand{\PSieteEntorno}{Linux 6.18.44-fc-v310, Python 3.13.18}


% Figura que está en tu Overleaf pero no en este paquete: si el archivo existe se
% incluye; si no, se dibuja un recuadro con el nombre que debe tener.
\newcommand{\figuraopcional}[2]{%
  \IfFileExists{#1}{\includegraphics[width=#2]{#1}}{%
    \fbox{\parbox[c][4cm][c]{#2}{\centering\footnotesize Subir a Overleaf:\\ \texttt{\detokenize{#1}}}}}}

% Nombres largos de código (con guion bajo) pueden cortarse al final de línea.
\let\eddieunderscore\_
\renewcommand{\_}{\eddieunderscore\allowbreak}
\setlength{\emergencystretch}{3em}
